Reasoning as before, one therefore has a canonical descent datum on the $\underline{\operatorname{Dyn}}(G_{S_i})$, allowing one to construct by descent a Dynkin $S$-scheme $\underline{\operatorname{Dyn}}(G)$.
\numpara{3.4}\label{XXIV.3.4}
This construction satisfies the following properties (which, moreover, essentially characterize it):
\begin{enumeratei}
\item To every reductive $S$-group is associated a Dynkin scheme $\underline{\operatorname{Dyn}}(G)$; to every isomorphism $u:G\isomto G'$ is functorially associated an isomorphism $\underline{\operatorname{Dyn}}(u):\underline{\operatorname{Dyn}}(G)\isomto\underline{\operatorname{Dyn}}(G')$.
\item If $S'$ is an $S$-scheme and $G$ a reductive $S$-group, one has
\[
\underline{\operatorname{Dyn}}(G\times_S S')\simeq\underline{\operatorname{Dyn}}(G)\times_S S'.
\]
\item If $\mathscr E$ is a pinning of $G$, defining the pinned root datum with Dynkin diagram $\Delta$, one has
\[
\underline{\operatorname{Dyn}}(G)\simeq\Delta_S.
\]
\item If $u$ is an inner automorphism of $G$, $\underline{\operatorname{Dyn}}(u)$ is the identity automorphism of $\underline{\operatorname{Dyn}}(G)$.
\end{enumeratei}
\numpara{3.5}\label{XXIV.3.5}
Let $S$ be a scheme and $G$ a reductive $S$-group. It is clear that the functor $\underline{\Aut}_{\mathrm{Dyn}}(\underline{\operatorname{Dyn}}(G))$ of automorphisms of $\underline{\operatorname{Dyn}}(G)$ for the Dynkin scheme structure is representable by a finite twisted constant $S$-scheme. By 3.4 (i) and (ii), one has a canonical morphism
\[
\underline{\Aut}_{S\text{-gr.}}(G)\to\underline{\Aut}_{\mathrm{Dyn}}(\underline{\operatorname{Dyn}}(G)),
\]
which, by (iv), factors through a morphism
\[
\underline{\operatorname{Autext}}(G)\to\underline{\Aut}_{\mathrm{Dyn}}(\underline{\operatorname{Dyn}}(G)).
\]
More generally, if $G$ and $G'$ are two reductive $S$-groups, one has a canonical morphism
\[
\underline{\operatorname{Isomext}}(G,G')\to\underline{\Isom}_{\mathrm{Dyn}}(\underline{\operatorname{Dyn}}(G),\underline{\operatorname{Dyn}}(G'));
\]
in particular, if $G'$ is an inner twisted form of $G$ (1.11), the Dynkin schemes $\underline{\operatorname{Dyn}}(G)$ and $\underline{\operatorname{Dyn}}(G')$ are isomorphic.
\numpara{3.6}\label{XXIV.3.6}
If $G$ is semisimple (resp. adjoint or simply connected), the morphism
\[
\underline{\operatorname{Autext}}(G)\to\underline{\Aut}_{\mathrm{Dyn}}(\underline{\operatorname{Dyn}}(G))
\]
is a monomorphism (resp. an isomorphism). Indeed, one may suppose $G$ pinned and one is reduced to the corresponding result for reduced pinned root data (cf. Exp.~XXI, 7.4.5).

One has an analogous result for the $\underline{\Isom}$; it follows in particular that two adjoint (resp. simply connected) semisimple $S$-groups are inner twisted forms of one another if and only if their Dynkin schemes are isomorphic.
\numpara{3.7}\label{XXIV.3.7}
One can give a different construction of the Dynkin scheme associated to a reductive group. Let $\mathscr R$ be a reduced pinned root datum and $G$ a reductive $S$-group of type $\mathscr R$; denote by $\Delta(\mathscr R)$ the Dynkin diagram defined by the root datum $\mathscr R$. One has (3.5) a canonical morphism
\[
A_S(\mathscr R)=\underline{\Aut}_{S\text{-gr.}}(\operatorname{\acute Ep}_S(\mathscr R))\to\underline{\Aut}_{\mathrm{Dyn}}(\Delta(\mathscr R)_S).
\]
The reductive $S$-group $G$ corresponds (1.17) to a bundle $\underline{\Isom}_{S\text{-gr.}}(\operatorname{\acute Ep}_S(\mathscr R),G)$, principal homogeneous under $A_S(\mathscr R)$. The associated bundle under $\underline{\Aut}_{\mathrm{Dyn}}(\Delta(\mathscr R)_S)$ corresponds to a form over $S$ of $\Delta(\mathscr R)_S$: it is $\underline{\operatorname{Dyn}}(G)$; in other words, this associated bundle is nothing other than $\underline{\Isom}_{\mathrm{Dyn}}(\Delta(\mathscr R)_S,\underline{\operatorname{Dyn}}(G))$. In this latter form the proof is immediate.
\numpara{3.8}\label{XXIV.3.8}\textbf{Dynkin scheme and Killing pairs.}
Let $S$ be a scheme, $G$ a reductive $S$-group, and $B\supset T$ a Killing pair of $G$. There exists a canonical morphism
\[
i:\underline{\operatorname{Dyn}}(G)\to\underline{\Hom}_{S\text{-gr.}}(T,\mathbf G_{\mathrm m,S})
\]
which identifies $\underline{\operatorname{Dyn}}(G)$ with the ``scheme of simple roots of $B$ relative to $T$''; this morphism is immediately defined by descent from the pinned case. Observe moreover that the datum of $T$ and $i$ allows one to reconstruct $B$ (``one-to-one correspondence between systems of simple roots and systems of positive roots'').

It follows from the preceding description of $\mathscr D=\underline{\operatorname{Dyn}}(G)$ that there exists a \emph{canonical root} of $B_{\mathscr D}$ relative to $T_{\mathscr D}$: this root $\alpha_{\mathscr D}$ is the image under $i(\mathscr D)$ of the identity morphism of $\mathscr D$. One deduces a \emph{canonical invertible $\OO_{\mathscr D}$-module $\mathfrak g^{\mathscr D}$}:
\[
\mathfrak g^{\mathscr D}=(\mathfrak g\otimes_{\OS}\OO_{\mathscr D})^{\alpha_{\mathscr D}}.
\]
In the pinned case one has
\[
\mathscr D=\coprod_{\alpha\in\Delta}S_\alpha,
\]
where each $S_\alpha$ is a copy of $S$, and $\mathfrak g^{\mathscr D}$ is then the $\OO_{\mathscr D}$-module which induces $\mathfrak g^\alpha$ on $S_\alpha$, for every $\alpha\in\Delta$.
\numpara{3.9}\label{XXIV.3.9}\textbf{Quasi-pinnings. Quasi-pinned groups.}
If $G$ is a reductive $S$-group, one calls a \emph{quasi-pinning} of $G$ the datum:
\begin{enumeratei}
\item of a Killing pair $B\supset T$ of $G$,
\item of a section $X\in\Gamma(\underline{\operatorname{Dyn}}(G),\mathfrak g^{\mathscr D})^\times$.
\end{enumeratei}
One says that a reductive $S$-group is \emph{quasi-splittable} if it possesses a quasi-pinning. One calls a \emph{quasi-pinned group} a reductive group equipped with a quasi-pinning.

Let $B\supset T$ be a Killing pair of the reductive $S$-group $G$; then $G$ admits a quasi-pinning relative to this Killing pair if and only if $\mathfrak g^{\mathscr D}$ has a section nonzero at every point, i.e. if the element of $\Pic(\underline{\operatorname{Dyn}}(G))$ defined by $\mathfrak g^{\mathscr D}$ is zero. Suppose in particular that $S$ is semi-local; then $\underline{\operatorname{Dyn}}(G)$ is also semi-local, hence $\Pic(\underline{\operatorname{Dyn}}(G))=0$. One deduces:
\begin{proposition}{3.9.1}\label{XXIV.3.9.1}
Let $S$ be a semi-local scheme and $G$ a reductive $S$-group. For $G$ to be quasi-splittable, it is necessary and sufficient that it possess a Borel subgroup.
\end{proposition}
\begin{proof}
\nde{We have given the reference to Exp.~XXII, 5.9.7 in detail.}Indeed, $S$ is affine, hence, by the first assertion of Exp.~XXII, 5.9.7, if $G$ possesses a Borel subgroup $B$, it also possesses a Killing pair $B\supset T$. Then, since $\Pic(\mathscr D)=0$, $\mathfrak g^{\mathscr D}$ has an everywhere nonzero section $X$.
\end{proof}
Let again $A$ be a semi-local ring and $S=\Spec(A)$; observe now that for every reductive $S$-group $G$ the morphism $\underline{\operatorname{Bor}}(G)\to S$ is surjective (since $G_{\bar s}$ possesses Borel subgroups, for every $s\in S$) and smooth and projective (Exp.~XXII, 5.8.3), hence has sections after a finite surjective étale extension of the base.\nde{The original referred to EGA~IV, §~24, which did not appear. One must use ``Bertini's theorem''. Let us give the argument in detail, as indicated to us by O.~Gabber. Let $X\to S$ be a surjective, smooth and projective morphism; replacing $S$ by a connected component, one may suppose that $X/S$ is everywhere of relative dimension $d$ (cf. EGA~IV$_4$, 17.10.2). One may also suppose that $X$ is a closed subscheme of a projective space $\mathbf P^n_S=\mathbf P(E)$, where $E$ is a free $A$-module of rank $n+1$ (cf. EGA~II, 5.3.3). Let $s_1,\ldots,s_r$ be the closed points of $S$; by Bertini's theorem (see for example [Jou83], I 6.10), there exists an open subset $U$ of the product $P=\mathbf P(E^*)^d$, with nonempty fibers, such that for every point $u=(f_1,\ldots,f_d)$ of $U_{s_i}$, the intersection of $X_{\kappa(u)}$ with the $d$ hyperplanes of $\mathbf P^n_{\kappa(u)}$ defined by $u$ is étale over $\kappa(u)$. After shrinking $U$, one may suppose that $U$ is the complement in affine space $\mathbf A^{nd}_S$ of the zero locus $\mathscr V(Q)$ of a certain polynomial $Q$ of degree $m>0$. One then sees easily (by induction on the number of variables) that $\mathscr V(Q)$ cannot contain all the rational points of $\mathbf A^{nd}_{s_i}$ if $|\kappa(s_i)|>m$. Since the morphism $A\to\prod_i\kappa(s_i)$ is surjective, one can find a monic polynomial $R\in A[X]$ of degree $m$ whose image in $\kappa(s_i)[X]$ is irreducible if $\kappa(s_i)$ is finite, and has $m$ distinct roots if $\kappa(s_i)$ is infinite. Put $A'=A[X]/(R)$ and $S'=\Spec(A')$; then $S'\to S$ is étale, finite and surjective, and the residue field at each of its closed points $s'_1,\ldots,s'_t$ has cardinality $\geq2^m$. Then $U$ has a rational point $u_i$ over each closed point of $S'$, and since $A'\to\prod_i\kappa(s_i)$ is surjective, these lift to a section $u$ of $P_{S'}$. Denote by $Z$ the intersection of $X_{S'}$ with the $d$ hyperplanes of $\mathbf P^n_{S'}$ defined by $u$, and by $V$ the open subset of $Z$ consisting of points where $Z$ is étale over $S'$. By EGA~IV$_3$, 11.3.8, $V$ contains the fibers $Z_{s'_i}$ for every $i$; since $\pi:Z\to S'$ is proper, it follows that the closed subset $\pi(Z-V)$ is empty, whence $V=Z$. Then $Z\to S'$ is surjective, étale and proper, hence finite, as is the composite $Z\to S$, and this supplies the desired section of $X\to S$.}
% typo?: the editorial note prints A' -> product kappa(s_i), without primes; kept as printed.
One deduces the
\begin{corollary}{3.9.2}\label{XXIV.3.9.2}
Let $S$ be a semi-local scheme and $G$ a reductive $S$-group. There exists an étale, finite and surjective morphism $S'\to S$ such that $G_{S'}$ is quasi-splittable.
\end{corollary}
\begin{remark}{3.9.3}\label{XXIV.3.9.3}
Under the preceding conditions, let $T$ be a maximal torus of $G$ (cf. Exp.~XIV, 3.20); then one may suppose in addition that $G_{S'}$ is quasi-splittable relative to $T_{S'}$: it suffices to apply the preceding reasoning to the ``scheme of Borel subgroups containing $T$'', which is finite and étale over $S$ (Exp.~XXII, 5.5.5 (ii)).
\end{remark}
\numpara{3.10}\label{XXIV.3.10}
Let $\mathscr E$ and $\mathscr E'$ be two quasi-pinnings of the reductive $S$-group $G$. There exists a \emph{unique} inner automorphism of $G$ taking $\mathscr E$ to $\mathscr E'$. Indeed, one is immediately reduced to the split case, where the assertion has already been proved (1.5; indeed, it suffices to observe that for an inner automorphism of $G$ it amounts to the same thing to respect a pinning or the underlying quasi-pinning). One concludes, as in no. 1, that a quasi-pinning of the reductive $S$-group $G$ defines a splitting $h$ of the exact sequence:
\[
\xymatrix@C=2em{1\ar[r]&\ad(G)\ar[r]&\underline{\Aut}_{S\text{-gr.}}(G)\ar@<-.6ex>[r]_p&\underline{\operatorname{Autext}}(G)\ar@<-.6ex>@{-->}[l]_h\ar[r]&1,}
\]
the image of $h$ being the subgroup of $\underline{\Aut}_{S\text{-gr.}}(G)$ which leaves the quasi-pinning invariant.

Likewise, if $G$ and $G'$ are two quasi-pinned $S$-groups, one naturally defines the subfunctor
\[
\underline{\Isom}_{S\text{-gr. q-pin.}}(G,G')
\]
of $\underline{\Isom}_{S\text{-gr.}}(G,G')$; the projection from $\underline{\Isom}_{S\text{-gr.}}(G,G')$ onto $\underline{\operatorname{Isomext}}(G,G')$ induces an isomorphism
\begin{equation}\tag{$*$}
\underline{\Isom}_{S\text{-gr. q-pin.}}(G,G')\isomto\underline{\operatorname{Isomext}}(G,G').
\end{equation}
\begin{theorem}{3.11}\label{XXIV.3.11}
Let $S$ be a scheme, $\mathscr R$ a reduced pinned root datum such that $\operatorname{\acute Ep}_S(\mathscr R)$ exists (cf. Exp.~XXV), and $E$ the group of its automorphisms. Consider the following three categories:
\begin{enumeratei}
\item The category $\underline{\mathrm{Rev}}$ of principal Galois coverings of $S$ with group $E$ (the morphisms are the isomorphisms of $E$-bundles).
\item The category $\underline{\mathrm{Redext}}$ whose objects are the reductive $S$-groups of type $\mathscr R$ (Exp.~XXII, 2.7), the morphisms from $G$ to $G'$ being the elements of $\underline{\operatorname{Isomext}}(G,G')(S)$.
\item The category $\underline{\mathrm{Q\acute ep}}$ of quasi-pinned reductive $S$-groups of type $\mathscr R$ (the morphisms are the isomorphisms respecting the quasi-pinnings).
\end{enumeratei}
These three categories are equivalent: more precisely, one has a diagram of functors, commutative up to functorial isomorphisms
\[
\xymatrix{\underline{\mathrm{Rev}}\ar[rr]^{\mathrm{q\acute ep}}&&\underline{\mathrm{Q\acute ep}}\ar[dl]^i\\
&\underline{\mathrm{Redext}}\ar[ul]^{\mathrm{rev}}&.}
\]
\end{theorem}
We shall describe these three functors below, leaving to the reader the task of verifying the commutativity of the diagram.\nde{In 3.11.4 we have sketched the verification that the functors $i$, $\mathrm{rev}$ and $\mathrm{q\acute ep}$ are fully faithful, and that the composite $\mathrm{rev}\circ i\circ\mathrm{q\acute ep}$ is isomorphic to the identity functor of $\underline{\mathrm{Rev}}$.}
\numpara{3.11.1}\label{XXIV.3.11.1}\textbf{The functor $i$.}
It is the evident functor: $i(G)=G$, $i(f)=$ image of $f$ under the isomorphism of 3.10 ($*$).
\numpara{3.11.2}\label{XXIV.3.11.2}\textbf{The functor $\mathrm{q\acute ep}$.}
Let $S'$ be a principal Galois covering of $S$ with group $E$. Let $(G,T,M,R,\Delta,(X_\alpha)_{\alpha\in\Delta})$ be a pinning of $G=\operatorname{\acute Ep}_S(\mathscr R)$ and let $(X'_\alpha)_{\alpha\in\Delta}$ be the corresponding pinning of $G'=G_{S'}=\operatorname{\acute Ep}_{S'}(\mathscr R)$. By Exp.~XXIII 5.5 bis, there exists a unique group morphism
\[
\theta:E\to\underline{\Aut}_{S\text{-gr.}}(\operatorname{\acute Ep}_S(\mathscr R))=A(\mathscr R)(S)
\]
such that, for every $h\in E$, $\theta(h)$ induces $\mathrm D_S(h)$ on $T$ and $\Lie(\theta(h))(X_\alpha)=X_{h(\alpha)}$ for every $\alpha\in\Delta$.
Taking account of the action of $E$ on $S'$, one thus obtains an action of $E$ on $G'$, compatible with the action of $E$ on $S'$ and respecting the canonical quasi-pinning of $G'$.\nde{We have expanded the original in the above, to show that the action of $E$ on $G'$ is obtained by combining the natural action of $E$ on $\operatorname{\acute Ep}_{S'}(\mathscr R)$ and the given action on $S'$. The canonical Killing pair of $G'$ is preserved by this action, as is the quasi-pinning given by the element $\widetilde X'$ of $\Gamma(\Delta_{S'},\Lie(G'/S')^{\Delta_{S'}})$, equal to $X'_\alpha$ on the copy of $S'$ indexed by $\alpha$ (indeed, since $E$ acts on $\Delta_{S'}$ by permuting the copies of $S'$, one has $h(\widetilde X')=\widetilde X'$ for every $h\in E$).}Since $S'\to S$ is covering for the (fpqc) topology, it is a morphism of effective descent for the fibered category of affine morphisms, and one denotes by
\[
\mathrm{q\acute ep}(S')=\operatorname{Q\text{-}\acute Ep}_{S'/S}(\mathscr R)
\]
the \emph{quasi-pinned} $S$-group obtained by Galois descent.\nde{See, for example, TDTE~I, p.~22, Example 1. One can also describe $\mathrm{q\acute ep}(S')$ as the twist of $G=\operatorname{\acute Ep}_S(\mathscr R)$ by the $E$-torsor $S'/S$, i.e. the (fpqc) quotient sheaf of $S'\times_S G$ by the right action of $E$ defined by $(s',g)\cdot h=(s'h,h^{-1}(g))$. Since $G$ is affine over $S$ and $E$ acts on $G$ by group automorphisms, this sheaf is representable by an $S$-group $G^\sharp$, which is a ``twisted'' form of $G$, and $\mathscr D^\sharp=\underline{\operatorname{Dyn}}(G^\sharp)$ is the twist of the constant Dynkin scheme $\Delta_S$ by the torsor $S'/S$. Since $E$ normalizes $B$ and $T$, one similarly obtains a pair $(B^\sharp,T^\sharp)$ which is a Killing pair of $G^\sharp$. Moreover, let $\mathfrak g=\Lie(G/S)$ and $\mathfrak g^\sharp=\Lie(G^\sharp/S)$; for every $U\to S$, the sections of $\mathfrak g^\sharp$ over $U$ are the $E$-equivariant $S$-morphisms $U\times_S S'\to\mathbf W(\mathfrak g)$. Since $\mathscr D^\sharp\times_S S'$ is $E$-isomorphic to $\Delta\times S'$, equipped with the action $(\alpha,s')\cdot h=(h^{-1}(\alpha),s'h)$, one obtains that the morphism given by $(\alpha,s')\mto X_\alpha$ is a section of $\mathfrak g^\sharp$ over $\mathscr D^\sharp$ which is a quasi-pinning, i.e. an everywhere nonzero section of $(\mathfrak g^\sharp\otimes_{\OS}\OO_{\mathscr D^\sharp})^{\mathscr D^\sharp}$ (cf. 3.8).}
\numpara{3.11.3}\label{XXIV.3.11.3}\textbf{The functor $\mathrm{rev}$.}
Let $G$ be a reductive $S$-group of type $\mathscr R$. One puts
\[
\mathrm{rev}(G)=\underline{\operatorname{Isomext}}(\operatorname{\acute Ep}_S(\mathscr R),G),
\]
it is a principal homogeneous bundle under $E_S$ (cf. 1.10 and 1.19), i.e. an object of $\underline{\mathrm{Rev}}$.
\numpara{3.11.4}\label{XXIV.3.11.4}
\nde{We have added this paragraph.}It is clear, by the isomorphism 3.10 ($*$), that the functor $i$ is fully faithful. Moreover, if $G_0$, $G$, $G'$ are reductive $S$-groups, the natural morphism
\[
\underline{\Isom}(G_0,G)\times_S\underline{\Isom}(G,G')\to\underline{\Isom}(G_0,G')
\]
induces a morphism
\begin{equation}\tag{$\star$}
\underline{\operatorname{Isomext}}(G_0,G)\times_S\underline{\operatorname{Isomext}}(G,G')\to\underline{\operatorname{Isomext}}(G_0,G')
\end{equation}
since, for every $S'\to S$, if $g_0\in G_0(S')$, $g\in G(S')$, $u\in\underline{\Isom}(G_0,G)(S')$, and $v\in\underline{\Isom}(G,G')(S')$, one has the equality $v\circ\mathrm{int}(g)\circ u\circ\mathrm{int}(g_0)=v\circ u\circ\mathrm{int}(u^{-1}(g)g_0)$. Taking $G_0=\operatorname{\acute Ep}_S(\mathscr R)$, one obtains a map
\[
\begin{aligned}
\underline{\operatorname{Isomext}}(G,G')(S)\to\underline{\Hom}_{\mathrm{Rev}}\bigl(&\underline{\operatorname{Isomext}}(\operatorname{\acute Ep}_S(\mathscr R),G),\\
&\underline{\operatorname{Isomext}}(\operatorname{\acute Ep}_S(\mathscr R),G')\bigr)
\end{aligned}
\]
and to show that it is a bijection, one may suppose, by (fpqc) descent, that $G$ and $G'$ are pinned, in which case it is evident.
Likewise, if $S_1$, $S_2$ are two objects of $\underline{\mathrm{Rev}}$ and if one denotes by $S_i\times^E\operatorname{\acute Ep}_S(\mathscr R)$ the $S$-group $\mathrm{q\acute ep}(S_i)$ obtained by twisting $\operatorname{\acute Ep}_S(\mathscr R)$ by the $E$-torsor $S_i$ (cf. N.D.E.~(21)), one has a natural map
\[
\begin{aligned}
\underline{\Hom}_{\mathrm{Rev}}(S_1,S_2)\to\underline{\Hom}_{\mathrm{Q\acute ep}}\bigl(&S_1\times^E\operatorname{\acute Ep}_S(\mathscr R),\\
&S_2\times^E\operatorname{\acute Ep}_S(\mathscr R)\bigr),
\end{aligned}
\]
and to show that it is a bijection, one may suppose, by (fpqc) descent, that $S_1$ and $S_2$ are trivial $E$-bundles, in which case it is evident.

Thus the three functors of the diagram are fully faithful. Moreover, for every object $S'$ of $\underline{\mathrm{Rev}}$ one has a natural morphism
\[
S'\to\underline{\operatorname{Isomext}}(\operatorname{\acute Ep}_S(\mathscr R),S'\times^E\operatorname{\acute Ep}_S(\mathscr R))
\]
and to see that it is an isomorphism, one may suppose, by (fpqc) descent, that $S'$ is a trivial $E$-bundle, in which case it is evident.

One therefore has an isomorphism of functors $\mathrm{Id}_{\underline{\mathrm{Rev}}}\to\mathrm{rev}\circ i\circ\mathrm{q\acute ep}$. Since the functor $\mathrm{rev}$ is fully faithful one therefore obtains, for every reductive $S$-group $H$ (not necessarily quasi-split), a bijection
\[
\underline{\operatorname{Isomext}}(\operatorname{Q\text{-}\acute Ep}_{\mathrm{rev}(H)/S}(\mathscr R),H)(S)\simeq\Aut_{\mathrm{Rev}}(\mathrm{rev}(H))
\]
functorial in $H$. In particular, the identity morphism of $\mathrm{rev}(H)$ corresponds to a ``canonical'' element $u_0$ of
\[
\underline{\operatorname{Isomext}}(\operatorname{Q\text{-}\acute Ep}_{\mathrm{rev}(H)/S}(\mathscr R),H)(S);
\]
moreover, every
\[
u\in\underline{\operatorname{Isomext}}(\operatorname{Q\text{-}\acute Ep}_{\mathrm{rev}(H)/S}(\mathscr R),H)(S)
\]
corresponds to an automorphism $\phi_u$ of $\mathrm{rev}(H)$, and $\mathrm{q\acute ep}(\phi_u)$ is the unique automorphism of the pinned $S$-group $\operatorname{Q\text{-}\acute Ep}_{\mathrm{rev}(H)/S}(\mathscr R)$ such that $u_0\circ\mathrm{q\acute ep}(\phi_u)=u$. We have thus proved Theorem 3.11.
% typo?: source says pinned S-group rather than quasi-pinned here; kept as printed.
Let us develop one of the corollaries of 3.11:
\begin{corollary}{3.12}\label{XXIV.3.12}
For every reductive $S$-group $G$, there exists a quasi-pinned $S$-group $G_{\text{q-pin.}}$ and an ``outer isomorphism'' $u\in\underline{\operatorname{Isomext}}(G_{\text{q-pin.}},G)(S)$. The pair $(G_{\text{q-pin.}},u)$ is unique up to a unique isomorphism.
\end{corollary}
Indeed, one may suppose $G$ of constant type $\mathscr R$, and one takes $G_{\text{q-pin.}}=\operatorname{Q\text{-}\acute Ep}_{\mathrm{rev}(G)/S}(\mathscr R)$.
\numpara{3.12.1}\label{XXIV.3.12.1}
Observe that the datum of the pair $(G_{\text{q-pin.}},u)$ allows one to define canonically the $S$-scheme (cf. 1.11)
\[
Q=\underline{\operatorname{Isomint}}(G_{\text{q-pin.}},G),
\]\nde{We write $\underline{\operatorname{Isomint}}(G_{\text{q-pin.}},G)$ instead of $\underline{\operatorname{Isomint}}_u(G_{\text{q-pin.}},G)$ (cf. 1.11), since the pair $(G_{\text{q-pin.}},u)$ is unique up to a unique isomorphism.}
which is a principal homogeneous bundle under $\ad(G_{\text{q-pin.}})$, and whose datum ``is equivalent'' to that of the inner twisted form $G$ of $G_{\text{q-pin.}}$. Moreover $Q$ is nothing other than the ``scheme of quasi-pinnings of $G$'', a definition which accounts for its structure of principal homogeneous bundle (on the left) under $\ad(G)$ (3.10)---compare with 1.20.
\begin{proposition}{3.13}\label{XXIV.3.13}
Let $S$ be a scheme, $G$ an adjoint (resp. simply connected) semisimple $S$-group, $B\supset T$ a Killing pair of $G$, and $\underline{\operatorname{Dyn}}(G)$ the Dynkin $S$-scheme of $G$.
